\section{Lezione 8} \hfill 11/03/2026
\subsection{Specie sporadiche}
Il termine sporadico non significa che sono una rarità, ma che di norma non sono in grado di formare popolamenti puri.\\
Questo è spesso dovuto a questioni allelopatiche o facilità di attacco di parassiti per gli individui più grossi.\\
Svolgono un ruolo ecologico unico all'interno dell'ecosistema forestale.\\
Molte di queste specie sono Rosacee, anche se non tutte. Per il contesto di pianura, occorre citare: l'Acero campestre (Acer campestre), ed il Ciliegio palustre (Prunus avium), legato alla presenza di acqua.\\
\subsubsection{Acero campestre - Acer campestre}
L'acero campestre è stato spesso usato (coltivato) in passato per maritare la vite (al fine di darle sostegno). È facile da potare ed è molto presente nelle campagne; risulta essere di scarso interesse forestale.
\subsubsection{Ciliegio selvatico - Prunus avium}
Nel piano collinare ci sono più specie sporadiche: la regina è il Ciliegio selvatico.\\
È possibile trovarlo a gruppetti di 4-5 piante: se la pianta è morta in modo traumatico, il ciliegio produce polloni radicali mediante ricaccio; quindi, quelle piante hanno tutte lo stesso patrimonio genetico.\\
È una pianta mesofila: da giovane può crescere abbastanza in ombra, ma poi richiede in fretta la luce per svilupparsi.\\
Richiede terreni fertili, umidi e freschi.\\
La sua produzione è molto importante per:
\begin{itemize}
    \item il suo legno molto apprezzato;
    \item il suo ruolo ecologico, come per tutte le specie sporadiche (produce frutti molto apprezzati dai selvatici).
\end{itemize}
\vspace{1cm}
\subsubsection{Sorbi}
Un'altra specie di piante sporadiche del piano collinare è relativo a tutta la categoria dei sorbi, come per esempio:
\begin{itemize}
    \item Ciavardello (Sorbus torminalis): è la specie europea di maggior valore, con un legno pallido e senza venature;
    \item Sorbo da sorbole (Sorbus domestica): si trova nel piano collinare (più appenninico che alpino);
    \item Sorbo montano (Sorbus aria): anche detto farinaccio, si trova a quote più elevate, e come il pippo ha le foglie bianche nella pagina inferiore;\\
    \item Sorbo degli uccellatori (Sorbus aucuparia): del piano montano e subalpino, non ha un elevato interesse per il legname, ma è importante dal punto di vista ecologico.
\end{itemize} 
il pero ed il melo selvatico (in minore parte). il pero ha un bel legno.\\
oggi, la chioma del melo è piatta , per poter riuscire a passare con i macchinari. \\
i peri erano impalcati a 2 metri d'altezza \\
queste specie vanno tenute di conto e tutelate (sia per la biodiversità che per il valore economico). vanno gestite con una selvicoltura d'albero (una selvicoltura di individuo). la selvicoltura di popolamento tende a creare una condizione media, con un assortimento medio. con la selvicoltura d'albero si vuole trarre il meglio dai pochi individui. \\
la gestione di popolamento era un problema. in passato era la gestione a ceduo: venivano favorite le specie con forte capacità pollonifera, che portano a buttare molto (popolamenti puri, difficile che ci siano più specie, con poche specie sporadiche). in molti boschi, le specie sporadiche sono scomparse, o sono presenti solo come rinnovazione. \\
la coltivazione con la selvicoltura d'albero è composta da 2-3 fasi:
\begin{itemize}
    \item fase di accestimento, rinnovazione: 
    \item fase di qualificazione,  nella quale la pianta forma il suo fusto, a piacimento del selvicoltore. tendenzialmente si vuole almeno 2 metri di toppo, netto dei rami. dev'essere raggiunta nel più breve tempo possibile;
    \item  fase di dimensionamento, nella quale la pianta accresce di diametro (preferibilmente nel toppo)
\end{itemize} 
la fase di rinnovazione la si trova in bosco (tendenzialmente non si piantano). la tutela degli individui avviene liberando la pianta dalla competizione (luce necessaria affinchè possa crescere e non subire l'influenza dalle piante vicine) . non si deve liberare subito , bisogna spingere la pianta verso l'alto fino che ottiene i 2 metri di toppo. avendo il prodotto, è necessario che fotosintetizzi molto , liberando la chioma. il resto del bosco viene gestito normalmente (sia questo ceduo o fustaia). \\
?, 37 min \\
\subsection{Platano - }
rinnova con difficoltà, molto usato in PP per la produzione di legna da ardere . sono due specie di albaro ed un ibrido ... l'ibrido venne creato in inghilterra 2-300 anni fa.\\
il legno ha molti usi, cap pollinare caulinare e ... veniva piantato in filari e ceduati oppure piantata come alto fusto negli incroci delle strade. veniva usato come alberature stradali, visto che si fa potare facilmente. \\
non si fa una grande selvicoltura, è importante nella selvicoltura urbana. \\
gestire popolamenti è imporatnte per il cancro colorato e  perchè se lo si trova indica che quel bosco non lo era in passato (specie segnale).\\
\subsection{tasso - }
è una pianta che in natura potrebbe essere presente dalla pianura alla montagna. è pressochè scomparsa per due motivi: tutta la pianta è velenosa (tranne il frutto), ed è la specie più sciafila che abbiamo. \\
sotto al tasso non cresce nulla: è stato sfavorito anche per questo.\\
ha un legno molto elastico ed è stato usato in passato per le frecce. è una delle specie che ha dato origine alle prime leggi forestali\\ importante conservarli 
\subsection{Piano collinare - regione avanalpica}
forte capacità pollonifera caulinare, fino a 30-40 anni. 
\subsubsection{Rovere - Quercus petrea}
La larghezza massima delle foglia è nella sua metà, ha un lungo picciolo. non ha la ghianda picciolata. la corteccia della rovere è finemente screpolata. le ramificazioni sono sottili, che portano la chioma verso l'alto. la rovere è una specie esigente e termofila (non sopporta gelate tardive), vuole suoli freschi e mai troppo umidi. è eliofila come la farnia, ha una vita come la farnia (?).\\
queste due specie, quando sono a contatto, si ibridano (rovere e farnia).\\
si trova dove le condizioni sono ottime .\\
legno di ottima qualità ed è un peccato ceduarlo . \\
è stata eliminata, per rimpiazzare il terreno per la coltivazione dell'uva. non hanno dominanza attualmente , ma si possono trovare nei boschi misti
\subsubsection{Cerro - Quercus cerris}
è una strana quecia, riconoscibile per tre aspetti: foglia stretta, oblunga e rugosa, la cupola delle ghiande ha della peluria, ha scanalature rosse nella corteccia. \\
dal punto di vista ecologico è una quercia relativamente sciafila.\\
si trova molto in appennino, dove va a contatto con il faggio (vicino al limite del bosco) e frequentemente nelle alpi (in situazioni più ombrose, versanti a Nord ,più freschi). il legno del cerro è nervoso , non viene usato come legname da opera, ma usato come legname da ardere. le ghiande sono amare, non apprezzate dai maiali. in appennino si aveva un bosco misto di ceduo comporto di roverella (fustaia) e cerro (ceduo)\\
la cerreta in passato veniva gestita a ceduo, dal dopo guerra si è fatta una trasformazione all'alto fusto  ?   \\
è stato usato in ferrovia per la produzione delle traversine ferroviarie 
\subsubsection{Roverella}
si riconosce per il dimorfismo delle foglie, che sono pubescenti nella faccia .... tiene le foglie secche nella pianta d'inverno e le perde quando deve farne di nuove. è la classica quercia delle condizioni difficili (sia substrati acidi che calcarei.): quando non cresce niente, nasce la roverella . \\
dato che cresce in stazioni difficili, ha fusti irregolari, con legno nervoso. ottiama legna da ardere e normalmente gestita a ceduo . \\
entra in gioco negli orno-ostrieti (carpino nero e orniello). 
\subsubsection{orno-ostrieti}
in substrati calcari , nelle condizioni peggiori (perchè nelle migliori ci sono...)\\
bosco a dominanza di orniello e carpino nero, ai quali si possono mescolare latifoglie nobili, specie sporadiche ... roverella (che potrebbe prendere il sopravvento). \\
formazioni abbondanti negli appennini  e nelle alpi... tranne dal ticino al po (perchè i suoli sono acidi).\\
il frassino da manna appare per primo è un ottimo gestore dell'acqua. il carpino nero è più esigente e normalmente compare in seguito e se le condizioni migliorano tende a prendere il sopravvento.\\
la fascia ideale degli orno-ostrieti è incastrata tra il bosco misto di latifoglie di sotto ed il bosco di faggio di sopra. ogni volta che le condizioni migliorano, ci sono le incursioni dei popolamento di sopra e di sotto. in natura questo popolamento è una mescolanza tra gli altri due popolamento.\\
le condizioni migliori sono stat...\\
accrescimenti ridotti, turni lunghi, cedui matricinati , turni tra i 20 e 30 anni, : la struttura era poverissima e le sporadiche sono sparite. \\
 l'avvento del gas ha dato respiro a questi popolamenti, portando al ritorno di altre specie. 